%%%%%%%%%%%%%%%%%%%%%%%%
%
% $Autor: Hemanth Jadiswami Prabhakaran $
% $Datum: 2026-09-25 $
% $Pfad: GitHub/MBIDA_Master_Thesis/Manual/Chapters/en/08Troubleshooting.tex $
% $Version: 1 $
%
% $Project: MBIDA Thesis - Manual $
%
%%%%%%%%%%%%%%%%%%%%%%%%


\chapter{Troubleshooting}
\label{chap:troubleshooting}

This chapter is organised by what you see. Find the symptom in Figure~\ref{fig:troubleshootingTree} and go to the section it names. Each section states the symptom, the cause and the solution. Solutions marked \emph{(operator)} need access to the repository or the Streamlit Community Cloud workspace.

\begin{figure}[H]
  \centering
  \resizebox{\linewidth}{!}{% Requires: \usetikzlibrary{positioning}
\begin{tikzpicture}[
	node distance=1cm and 1cm, % vertical distance = 1cm
	sym/.style={man/boxorange, text width=4.1cm, font=\sffamily\scriptsize, align=left,
		minimum height=1.05cm, fill=white},
	grp/.style={man/decision, text width=2.0cm, font=\sffamily\scriptsize, align=center},
	]
	\node[man/terminal, font=\sffamily\small] (root) at (0,0) {What do you see?};
	
	% Group headers (1cm vertical gap below root level)
	\node[grp] (g2) [below=1cm of root] {An error or\\warning box};
	\node[grp] (g1) [left=3cm of g2]     {The page does\\not load};
	\node[grp] (g3) [right=3cm of g2]    {The chart looks\\wrong};
	
	\foreach \g in {g1,g2,g3} \draw[man/arrow] (root) -- (\g);
	
	% Column 1 (each node 1cm below previous border)
	\node[sym] (a1) [below=1cm of g1] {``This app has gone to sleep''\\[1pt]\textbf{$\rightarrow$ Section~\ref{sec:tsSleep}}};
	\node[sym] (a2) [below=1cm of a1] {Local: \texttt{streamlit: command not found}\\[1pt]\textbf{$\rightarrow$ Section~\ref{sec:tsCommand}}};
	\node[sym] (a3) [below=1cm of a2] {Local: \texttt{ModuleNotFoundError}\\[1pt]\textbf{$\rightarrow$ Section~\ref{sec:tsModule}}};
	\node[sym] (a4) [below=1cm of a3] {Local: port 8501 already in use\\[1pt]\textbf{$\rightarrow$ Section~\ref{sec:tsPort}}};
	
	% Column 2
	\node[sym] (b1) [below=1cm of g2] {Red box: ``Couldn't find one of the required parquet files''\\[1pt]\textbf{$\rightarrow$ Section~\ref{sec:tsParquet}}};
	\node[sym] (b2) [below=1cm of b1] {Yellow box: ``No data found for \dots''\\[1pt]\textbf{$\rightarrow$ Section~\ref{sec:tsNoData}}};
	\node[sym] (b3) [below=1cm of b2] {Node line ends with ``No column found for \dots''\\[1pt]\textbf{$\rightarrow$ Section~\ref{sec:tsNoColumn}}};
	
	% Column 3
	\node[sym] (c1) [below=1cm of g3] {No reconciled line, only actual + base\\[1pt]\textbf{$\rightarrow$ Section~\ref{sec:tsNoRecLine}}};
	\node[sym] (c2) [below=1cm of c1] {Colours hard to read / theme not followed\\[1pt]\textbf{$\rightarrow$ Section~\ref{sec:tsTheme}}};
	\node[sym] (c3) [below=1cm of c2] {A dropdown offers only ``(All)''\\[1pt]\textbf{$\rightarrow$ Section~\ref{sec:tsLocked}}};
	\node[sym] (c4) [below=1cm of c3] {Forecast line stops, or looks flat\\[1pt]\textbf{$\rightarrow$ Section~\ref{sec:tsFlat}}};
	
	% Downward arrows (all exactly 1cm)
	\draw[man/arrow] (g1) -- (a1); \draw[man/arrow] (a1) -- (a2);
	\draw[man/arrow] (a2) -- (a3); \draw[man/arrow] (a3) -- (a4);
	\draw[man/arrow] (g2) -- (b1); \draw[man/arrow] (b1) -- (b2); \draw[man/arrow] (b2) -- (b3);
	\draw[man/arrow] (g3) -- (c1); \draw[man/arrow] (c1) -- (c2);
	\draw[man/arrow] (c2) -- (c3); \draw[man/arrow] (c3) -- (c4);
	
	\node[man/label, below=1cm of b3] {If nothing matches:\\Section~\ref{sec:tsGeneral}};
\end{tikzpicture}}
  \caption{Finding the right section from the symptom}
  \label{fig:troubleshootingTree}
\end{figure}

% ------------------------------------------------------------------
\section{The Page Does Not Load}
\label{sec:tsPageGroup}

\subsection{``This app has gone to sleep''}
\label{sec:tsSleep}

\begin{description}[style=nextline,leftmargin=0.5cm,font=\normalfont\bfseries]
  \item[Symptom] Instead of the dashboard, a page says the app has gone to sleep due to inactivity.
  \item[Cause] Streamlit Community Cloud stops apps that have not been used for a while. This is normal behaviour of the free hosting.
  \item[Solution] Click \UI{Yes, get this app back up!} and wait up to a minute (Section~\ref{sec:instWake}). If the app does not start within five minutes: \emph{(operator)} open the logs and reboot the app (Section~\ref{sec:mtCloud}).
\end{description}

\subsection{The Page Stays Blank or Shows ``Oh no''}
\label{sec:tsBlank}

\begin{description}[style=nextline,leftmargin=0.5cm,font=\normalfont\bfseries]
  \item[Symptom] The page remains white, keeps loading, or Streamlit shows ``Oh no. Error running app.''
  \item[Cause] A slow or interrupted connection, a browser extension blocking scripts, or an error during the app's (re)start.
  \item[Solution] Reload the page (\keys{F5}); try a private browser window, which disables most extensions; try another browser. If the error persists for all users: \emph{(operator)} check the logs -- a failed package installation after a change is the most frequent reason -- and roll back the last change if needed (Section~\ref{sec:mtBackup}).
\end{description}

\subsection{\texttt{streamlit: command not found}}
\label{sec:tsCommand}

\begin{description}[style=nextline,leftmargin=0.5cm,font=\normalfont\bfseries]
  \item[Symptom] Local start fails with \texttt{command not found} (macOS/Linux) or ``is not recognized as the name of a cmdlet'' (Windows).
  \item[Cause] The virtual environment is not activated, or the packages were not installed into it.
  \item[Solution] Activate the environment (the prompt must start with \texttt{(.venv)}), then install and start again:
\begin{lstlisting}[style=shell]
source .venv/bin/activate          # Windows: .venv\Scripts\Activate.ps1
pip install -r Code/app/requirements.txt
python -m streamlit run Code/app/streamlit_app.py
\end{lstlisting}
  On Windows, if activation is refused because of the execution policy, run once: \texttt{Set-ExecutionPolicy -Scope CurrentUser RemoteSigned}.
\end{description}

\subsection{\texttt{ModuleNotFoundError}}
\label{sec:tsModule}

\begin{description}[style=nextline,leftmargin=0.5cm,font=\normalfont\bfseries]
  \item[Symptom] The browser or terminal shows \texttt{ModuleNotFoundError: No module named '\dots'}.
  \item[Cause and solution] Depends on the module name:
  \begin{itemize}
    \item \texttt{streamlit\_theme} -- the wrong package \texttt{streamlit-theme} was installed instead of \texttt{st-theme}. Fix: \texttt{pip uninstall streamlit-theme} and \texttt{pip install st-theme==1.2.3}.
    \item \texttt{config} -- the file \texttt{Code/config.py} is missing or the app file was moved out of \texttt{Code/app/}. Fix: restore the original folder layout (\texttt{git status} shows what changed).
    \item \texttt{pyarrow} or any other package -- the requirements were not installed completely. Fix: \texttt{pip install -r Code/app/requirements.txt}.
  \end{itemize}
\end{description}

\subsection{Port 8501 Already in Use}
\label{sec:tsPort}

\begin{description}[style=nextline,leftmargin=0.5cm,font=\normalfont\bfseries]
  \item[Symptom] Streamlit reports that port 8501 is not available, or it starts on 8502 instead.
  \item[Cause] Another Streamlit app -- often an earlier start of this dashboard -- is still running.
  \item[Solution] Stop the other app with \keys{Ctrl+C} in its terminal, or start this one on another port:
\begin{lstlisting}[style=shell]
streamlit run Code/app/streamlit_app.py --server.port 8600
\end{lstlisting}
\end{description}

% ------------------------------------------------------------------
\section{Error and Warning Messages}
\label{sec:tsMessageGroup}

\subsection{``Couldn't find one of the required parquet files''}
\label{sec:tsParquet}

\begin{description}[style=nextline,leftmargin=0.5cm,font=\normalfont\bfseries]
  \item[Symptom] A red box with this text and the name of the missing file; no chart (Figure~\ref{fig:tsParquet}).
  \item[Cause] One of the four result files is missing from \texttt{Code/artifacts/dashboard\_data/}, e.g.\ because it was deleted, not committed or not downloaded.
  \item[Solution] \emph{(operator)} Restore the files from Git or regenerate them (Section~\ref{sec:mtRegenerate}).
\end{description}

\begin{figure}[H]
  \centering
  \Screenshot[0.85\linewidth]{Images/08Troubleshooting/ErrorMissingParquet.png}
  \caption{Error message when a result file is missing}
  \label{fig:tsParquet}
\end{figure}

\subsection{``No data found for \dots''}
\label{sec:tsNoData}

\begin{description}[style=nextline,leftmargin=0.5cm,font=\normalfont\bfseries]
  \item[Symptom] A yellow box instead of the chart, naming a node such as \VAL{Total/CA/CA\_1/\dots}.
  \item[Cause] The dropdown list and the sales file do not match -- normally only possible if the result files come from different pipeline runs.
  \item[Solution] Choose another node to continue. \emph{(operator)} Make sure all four files were written by the same pipeline run; restore them from Git or regenerate them together (Section~\ref{sec:mtRegenerate}).
\end{description}

\subsection{``No column found for \dots''}
\label{sec:tsNoColumn}

\begin{description}[style=nextline,leftmargin=0.5cm,font=\normalfont\bfseries]
  \item[Symptom] The node line ends with \texttt{No column found for: MinTrace (wls\_struct)} (or another method), and that method's line is missing.
  \item[Cause] The technical name of the method in \texttt{Code/config.py} does not match the result files -- usually after the pipeline or the settings were changed.
  \item[Solution] \emph{(operator)} Align the method names in \texttt{config.py} with the result files (Section~\ref{sec:mtConfig}), or undo the last change.
\end{description}

% ------------------------------------------------------------------
\section{The Chart Looks Wrong}
\label{sec:tsChartGroup}

\subsection{No Reconciled Line}
\label{sec:tsNoRecLine}

\begin{description}[style=nextline,leftmargin=0.5cm,font=\normalfont\bfseries]
  \item[Symptom] Only the actual and the dotted base line are visible.
  \item[Cause] The \UI{Reconciliation method(s)} box is empty, or the line was hidden by clicking its legend entry.
  \item[Solution] Select at least one method in the box; click the greyed-out legend entry to show the line again.
\end{description}

\subsection{A Reconciled Line Hides Another Line}
\label{sec:tsOverlap}

\begin{description}[style=nextline,leftmargin=0.5cm,font=\normalfont\bfseries]
  \item[Symptom] Fewer lines are visible than the legend lists.
  \item[Cause] Two forecasts are identical and lie exactly on top of each other -- for example Top-Down and the base forecast at the Total, or Bottom-Up and the base forecast at item level (Section~\ref{sec:fnRec}).
  \item[Solution] None needed. Hide one line via the legend to check.
\end{description}

\subsection{Colours Hard to Read}
\label{sec:tsTheme}

\begin{description}[style=nextline,leftmargin=0.5cm,font=\normalfont\bfseries]
  \item[Symptom] Lines or labels have too little contrast, e.g.\ after switching between light and dark mode.
  \item[Cause] The chart detects the theme when it is drawn; in rare cases the first drawing happens before the theme is known.
  \item[Solution] Change any control once or reload the page (\keys{F5}); the chart is then drawn with the right colours. Choose \VAL{Light} in the page menu for printing.
\end{description}

\subsection{A Dropdown Offers Only ``(All)''}
\label{sec:tsLocked}

\begin{description}[style=nextline,leftmargin=0.5cm,font=\normalfont\bfseries]
  \item[Symptom] For example, \UI{Department} cannot be opened to any value.
  \item[Cause] A dropdown above it is still on \VAL{(All)}. This is intended (Section~\ref{sec:fnNavigate}).
  \item[Solution] Choose values in all dropdowns above it first, from top to bottom.
\end{description}

\subsection{Forecast Looks Flat, or Is Zero}
\label{sec:tsFlat}

\begin{description}[style=nextline,leftmargin=0.5cm,font=\normalfont\bfseries]
  \item[Symptom] The forecast is a horizontal line, identical for all three base models, or lies at zero.
  \item[Cause] \VAL{histavg} is flat by definition. For items that sell rarely, \VAL{ets} and \VAL{theta} fall back to \VAL{histavg}, and forecasts for items that hardly sell at all are close to zero (Section~\ref{sec:fnBase}).
  \item[Solution] None needed; this is correct behaviour. Compare with the department or store above the item to see the models' differences.
\end{description}

% ------------------------------------------------------------------
\section{General Procedure}
\label{sec:tsGeneral}

If none of the symptoms above matches, work through these steps in order:
\begin{enumerate}
  \item Reload the page (\keys{F5}); try a private window and a second browser.
  \item Check whether the problem also occurs online \emph{and} locally. If only online: \emph{(operator)} check the logs and reboot (Section~\ref{sec:mtCloud}). If only locally: reset the local installation (Section~\ref{sec:mtLocalReset}).
  \item Compare with the verification checklist (Section~\ref{sec:mtVerify}) to narrow down which function fails.
  \item Report the problem as described in Section~\ref{sec:helpReport}.
\end{enumerate}
